\section*{Chapitre \ref{ch:dv:barriers-and-digitals} --- Barrières et digitales}

\begin{solution}{exo:dv:barriers-and-digitals:1}
Exactement l'un des deux verse 1 à l'échéance, quoi qu'il arrive, si bien qu'ensemble ils versent 1 avec certitude~: une obligation zéro-coupon, $e^{-rT}$.
Leurs deltas ont donc une somme nulle~: le delta de l'\omterm{def:dv:barriers-and-digitals:digital}{option digitale} de vente est l'opposé de celui du call.
\end{solution}

\begin{solution}{exo:dv:barriers-and-digitals:2}
L'\omterm{def:dv:barriers-and-digitals:digital}{option digitale} vaut $-\partial C/\partial K$. Avec un skew de volatilité à la baisse, la volatilité du call diminue quand le strike augmente, si bien que le call
perd de la valeur avec le strike plus vite qu'à volatilité constante, et $-\partial C/\partial K$ gagne le terme
$-\mathcal V\,\partial\sigma/\partial K>0$. En termes de probabilités~: la densité avec skew a une queue gauche épaisse et, pour garder
le forward fixe, davantage de masse juste au-dessus du forward, si bien que la probabilité de finir au-dessus d'un strike à la monnaie est plus élevée
que dans le modèle à volatilité plate.
\end{solution}

\begin{solution}{exo:dv:barriers-and-digitals:3}
La parité in-out n'est pas respectée~: $7.23+1.70=8.93\ne8.83$. Avec des rebates nuls, les deux doivent totaliser le prix de l'option vanille~; le knock-in
devrait valoir 1.60, à moins qu'un des prix n'utilise des paramètres différents (barrière, observation ou volatilité).
\end{solution}

\begin{solution}{exo:dv:barriers-and-digitals:4}
Avec $K/H=1.11$, les puts de strike $H^2/K=81$ valent 1.498, et le call down-and-in vaut aussi 1.498. Si la barrière n'est jamais touchée
l'indice reste au-dessus de 90 et clôture au-dessus de 90, donc au-dessus du strike des puts, qui est de 81~: les puts expirent sans valeur, comme le
knock-in.
\end{solution}

\begin{solution}{exo:dv:barriers-and-digitals:5}
Le facteur est $e^{0.5826\times0.2\times\sqrt{1/252}}=e^{0.0073}$~: la barrière à 90 descend à 89.34 et la barrière à 120
monte à 120.88, en s'éloignant du spot dans les deux cas.
\end{solution}

\begin{solution}{exo:dv:barriers-and-digitals:6}
Quand le spot monte vers 120, la valeur intrinsèque du call augmente, mais la probabilité de toucher la barrière et
de tout perdre~; près de la barrière, le second effet domine, si bien que la valeur baisse quand le spot monte et que le delta est
négatif. La valeur atteint son maximum là où les deux effets se compensent, bien en dessous de la barrière (vers 111, avec un mois restant, où elle vaut 7.68).
\end{solution}

\begin{solution}{exo:dv:barriers-and-digitals:7}
Continu~: 0.162. Quotidien~: Monte Carlo 0.205, décalé 0.207~; hebdomadaire 0.262 et 0.270~; mensuel 0.383 et 0.430~; quatre dates 0.564
et 0.735. Le décalage est efficace au quotidien, et se dégrade plus vite que pour le call quand l'observation se raréfie, parce que le
payoff du put est fort à la barrière ($K-H=10$), là où le développement est le moins précis.
\end{solution}

\begin{solution}{exo:dv:barriers-and-digitals:8}
Un ordre stop-loss placé à la barrière est exécuté au prochain cours disponible. Si le marché saute le niveau (une
annonce, un changement de politique, un mouvement de nuit), l'exécution a lieu bien au-delà du niveau et le P\&L de la couverture est égal au gap multiplié par la
position, comme le 15 janvier 2015. Le \omterm{def:dv:barriers-and-digitals:gap}{risque de gap} doit être valorisé (modèles à sauts) et limité (par l'exposition à chaque niveau de barrière),
et non ignoré.
\end{solution}

\begin{solution}{pb:dv:barriers-and-digitals:1}
\textbf{1.} 4.84 millions à 16.4~\%~; 5.77 millions avec le smile.
\textbf{2.} 0.93 million~: $-\mathcal V\,\partial\sigma/\partial K$, positif parce que le skew fait baisser la volatilité avec le
strike.
\textbf{3.} 1.25 million par point d'indice, soit 6.25 milliards d'indice.
\textbf{4.} Il faudrait l'acheter ou la vendre en une minute et la renverser à chaque tick en franchissant le strike~: bien au-delà
de la profondeur du marché, et la couverture perdrait à chaque fois la fourchette.
\textbf{5.} Pour le détenteur, élevé et positif juste sous le strike et élevé et négatif juste au-dessus~; le vendeur a
l'inverse.
\textbf{6.} Achats de calls en $K-\varepsilon$, ventes de calls en $K$, pour une taille $N/\varepsilon$~: il verse au moins autant que l'\omterm{def:dv:barriers-and-digitals:digital}{option
digitale} partout, si bien que le vendeur qui la détient n'est jamais à découvert à l'échéance.
\textbf{7.} $\varepsilon=10\,000\,000/50\,000=200$ points d'indice~: le payoff du call spread varie d'au plus 50\,000 par point.
\textbf{8.} 6.66 millions~; la surcouverture vaut 0.89 million en sus des 5.77 millions de la digitale.
\textbf{9.} 0.25, 0.48 et 1.56 millions.
\textbf{10.} 5 millions, soit la moitié du notionnel de la digitale, au lieu de rien~: la surcouverture paie pour cette zone.
\textbf{11.} 7.23 et 1.60.
\textbf{12.} 7.48 (erreur type 0.03) par simulation et 7.45 par décalage.
\textbf{13.} 6.15 contre 6.51~: la \omterm{def:dv:local-volatility:model}{volatilité locale} augmente quand l'indice baisse vers la barrière, si bien que la barrière est touchée plus
souvent qu'avec une volatilité plate.
\textbf{14.} 0.141.
\textbf{15.} 0.592 au moment du contact, 0.581 à l'échéance~: payer au contact évite l'actualisation entre le contact et l'échéance.
\textbf{16.} Le client achète un payoff qui domine la digitale~; le surcoût paie le payoff supplémentaire entre
4\,800 et 5\,000, et pour une couverture que le desk peut réellement gérer.
\textbf{17.} Quand le desk peut accepter un risque de pin en échange d'un prix proche de celui de la digitale, par exemple quand il
peut le compenser avec d'autres positions aux strikes proches~; un spread centré est moins cher mais peut laisser le vendeur à découvert
à l'échéance.
\textbf{18.} Une \omterm{def:dv:barriers-and-digitals:static}{couverture statique} lorsque c'est possible (un calendrier d'options sur la barrière), avec une limite sur le delta à la barrière et une
réserve pour les gaps~; la couverture dynamique seule échoue précisément quand la barrière est franchie par un saut.
\textbf{19.} 200 points d'indice~: un call spread de 50\,000 par point, de 4\,800 à 5\,000, qui facture au client une
surcouverture de 0.89 million sur la digitale de 10 millions (5.77 millions).
\textbf{20.} Une discontinuité, dans le payoff ou à la barrière, où le delta devient non borné et où la couverture doit être
remplacée par une surcouverture, un portefeuille statique ou une réserve.
\end{solution}

\begin{solution}{iq:dv:barriers-and-digitals:1}
Évaluer la digitale comme l'opposé de la dérivée du prix du call par rapport au strike, ce qui fait intervenir la pente du smile. La couvrir comme un call spread~:
enregistrer la digitale vendue au prix du spread dominant (la surcouverture) et couvrir en delta le spread, dont le delta est
borné.
\iqlookfor{le terme de smile et la surcouverture.}
\end{solution}

\begin{solution}{iq:dv:barriers-and-digitals:2}
La densité absorbée du log-prix à une barrière basse est la densité libre moins son symétrique par rapport à la barrière, pondérée
par le facteur de dérive. En intégrant le payoff du call par rapport à elle, on obtient
$C_{\mathrm{BS}}(S)-(H/S)^{2\nu/\sigma^2}C_{\mathrm{BS}}(H^2/S)$ pour $K\ge H$~; vérifier qu'elle résout l'\omterm{def:dv:black-scholes-three-ways:pde}{équation de Black--Scholes}, qu'elle s'annule
sur la barrière et a le bon payoff.
\iqlookfor{l'image, et la vérification aux bords.}
\end{solution}

\begin{solution}{iq:dv:barriers-and-digitals:3}
Une \omterm{def:dv:barriers-and-digitals:digital}{option digitale} vaut $-\partial C/\partial K$, elle dépend donc de la pente du smile au strike~: $-\mathcal V\partial\sigma/\partial K$.
Avec un skew à la baisse, ce terme est positif pour une digitale d'achat (qui vaut plus que son prix de Black--Scholes) et négatif pour une
digitale de vente.
\iqlookfor{la dérivée et le signe.}
\end{solution}

\begin{solution}{iq:dv:barriers-and-digitals:4}
Non~: un knock-out à observation quotidienne ne voit pas les franchissements entre deux observations, si bien qu'il vaut plus que la formule continue. Comparer
avec la formule où la barrière est décalée de $e^{\beta\sigma\sqrt{\Delta t}}$ en s'éloignant du spot~; si les deux concordent à l'intérieur de
l'erreur de simulation, le Monte Carlo est correct.
\iqlookfor{observation discrète face à l'observation continue, et le décalage.}
\end{solution}

\begin{solution}{iq:dv:barriers-and-digitals:5}
La perte survient si le sous-jacent saute par-dessus une barrière ou un strike où la valeur de la position ou la couverture change brusquement. La mesurer
par scénarios~: le P\&L du portefeuille pour des sauts de plusieurs tailles de part et d'autre de chaque niveau de barrière, en supposant les couvertures exécutées au
cours obtenu après le saut~; limiter l'exposition concentrée sur un même niveau, et la valoriser avec des modèles à sauts ou des réserves.
\iqlookfor{le scénario, la concentration sur les niveaux et la valorisation.}
\end{solution}

\begin{solution}{iq:dv:barriers-and-digitals:6}
Avec un carry nul et un smile symétrique, détenir $K/H$ puts de strike $H^2/K$~: sur la barrière, ils valent exactement le call, si bien
qu'au moment du contact, on les échange contre le call~; si la barrière n'est jamais touchée, ils expirent sans valeur. La méthode échoue avec un carry (la symétrie suppose un forward sans dérive),
avec un skew (les puts situés sous la barrière sont valorisés à une volatilité plus élevée que le call situé au-dessus), et avec des gaps à travers la
barrière.
\iqlookfor{la symétrie, et ses trois modes d'échec.}
\end{solution}
